\section{AHSME 1998}
        \begin{enumerate}
        	\item (\textit{AHSME 1998}) Each of the sides of five congruent rectangles is labeled with an integer. In rectangle A, $w = 4, x = 1, y = 6, z = 9$. In rectangle B, $w = 1, x = 0, y = 3, z = 6$. In rectangle C, $w = 3, x = 8, y = 5, z = 2$. In rectangle D, $w = 7, x = 5, y = 4, z = 8$. In rectangle E, $w = 9, x = 2, y = 7, z = 0$. These five rectangles are placed, without rotating or reflecting, in position as below. Which of the rectangle is the top leftmost one?


 
                    \begin{figure}[H]
                        \centering
                        \includegraphics[width=0.5\textwidth]{Problem/AHSME_I_1998_1_Problem_Figure_1.png}
                    \end{figure}
                    

 $\mathrm{(A)\ } A \qquad \mathrm{(B) \ }B \qquad \mathrm{(C) \  } C \qquad \mathrm{(D) \  } D \qquad \mathrm{(E) \  }E$


 

	\item (\textit{AHSME 1998}) Letters $A,B,C,$ and $D$ represent four different digits selected from $0,1,2,\ldots ,9.$ If $(A+B)/(C+D)$ is an integer that is as large as possible, what is the value of $A+B$?


 $\mathrm{(A) \  }13 \qquad \mathrm{(B) \  }14 \qquad \mathrm{(C) \  } 15\qquad \mathrm{(D) \  }16 \qquad \mathrm{(E) \  } 17$


 

	\item (\textit{AHSME 1998}) If $\texttt{a,b,}$ and $\texttt{c}$ are digits for which


 
 \[\begin{array}{rccc}&\ \texttt{7}& \texttt{a}&\texttt{2}\\ -&\ \texttt{4}&\texttt{8}&\texttt{b} \\  \hline  &\ \texttt{c}&\texttt{7}& \texttt{3} \end{array}\]


 then $\texttt{a+b+c =}$


 $\mathrm{(A) \  }14 \qquad \mathrm{(B) \  }15 \qquad \mathrm{(C) \  }16 \qquad \mathrm{(D) \  }17 \qquad \mathrm{(E) \  }18$


 

	\item (\textit{AHSME 1998}) Define $[a,b,c]$ to mean $\frac {a+b}c$, where $c \neq 0$. What is the value of 


 $\left[[60,30,90],[2,1,3],[10,5,15]\right]?$

 $\mathrm{(A) \ }0 \qquad \mathrm{(B) \ }0.5 \qquad \mathrm{(C) \ }1 \qquad \mathrm{(D) \ }1.5 \qquad \mathrm{(E) \ }2$


 

	\item (\textit{AHSME 1998}) If $2^{1998}-2^{1997}-2^{1996}+2^{1995} = k \cdot 2^{1995},$ what is the value of $k$?


 $\mathrm{(A) \ } 1 \qquad \mathrm{(B) \ } 2 \qquad \mathrm{(C) \ } 3 \qquad \mathrm{(D) \ } 4 \qquad \mathrm{(E) \ } 5$


 

	\item (\textit{AHSME 1998}) If $1998$ is written as a product of two positive integers whose difference is as small as possible, then the difference is 


 $\mathrm{(A) \ }8 \qquad \mathrm{(B) \ }15 \qquad \mathrm{(C) \ }17 \qquad \mathrm{(D) \ }47 \qquad \mathrm{(E) \ } 93$


 

	\item (\textit{AHSME 1998}) If $N > 1$, then $\sqrt[3]{N\sqrt[3]{N\sqrt[3]{N}}} =$


 $\mathrm{(A) \ } N^{\frac 1{27}} \qquad \mathrm{(B) \ } N^{\frac 1{9}} \qquad \mathrm{(C) \ } N^{\frac 1{3}} \qquad \mathrm{(D) \ } N^{\frac {13}{27}} \qquad \mathrm{(E) \ } N$


 

	\item (\textit{AHSME 1998}) A square with sides of length $1$ is divided into two congruent trapezoids and a pentagon, which have equal areas, by joining the center of the square with points on three of the sides, as shown. Find $x$, the length of the longer parallel side of each trapezoid. 


 
                    \begin{figure}[H]
                        \centering
                        \includegraphics[width=0.5\textwidth]{Problem/AHSME_I_1998_8_Problem_Figure_1.png}
                    \end{figure}
                    

 $\mathrm{(A) \ } \frac 35 \qquad \mathrm{(B) \ } \frac 23 \qquad \mathrm{(C) \ } \frac 34 \qquad \mathrm{(D) \ } \frac 56 \qquad \mathrm{(E) \ } \frac  78$


 

	\item (\textit{AHSME 1998}) A speaker talked for sixty minutes to a full auditorium. Twenty percent of the audience heard the entire talk and ten percent slept through the entire talk. Half of the remainder heard one third of the talk and the other half heard two thirds of the talk. What was the average number of minutes of the talk heard by members of the audience?


 $\mathrm{(A) \ } 24 \qquad \mathrm{(B) \ } 27\qquad \mathrm{(C) \ }30 \qquad \mathrm{(D) \ }33 \qquad \mathrm{(E) \ }36$


 

	\item (\textit{AHSME 1998}) A large square is divided into a small square surrounded by four congruent rectangles as shown. The perimeter of each of the congruent rectangles is $14$. What is the area of the large square?


 
                    \begin{figure}[H]
                        \centering
                        \includegraphics[width=0.5\textwidth]{Problem/AHSME_I_1998_10_Problem_Figure_1.png}
                    \end{figure}
                    

 $\mathrm{(A) \ }49 \qquad \mathrm{(B) \ }64 \qquad \mathrm{(C) \ }100 \qquad \mathrm{(D) \ }121 \qquad \mathrm{(E) \ }196$


 

	\item (\textit{AHSME 1998}) Let $R$ be a rectangle. How many circles in the plane of $R$ have a diameter both of whose endpoints are vertices of $R$?


 $\mathrm{(A) \ }1 \qquad \mathrm{(B) \ }2 \qquad \mathrm{(C) \ }4 \qquad \mathrm{(D) \ }5 \qquad \mathrm{(E) \ }6$


 

	\item (\textit{AHSME 1998}) How many different prime numbers are factors of $N$ if


 $\log_2 ( \log_3 ( \log_5 (\log_ 7 N))) = 11?$

 $\mathrm{(A) \ }1 \qquad \mathrm{(B) \ }2 \qquad \mathrm{(C) \ }3 \qquad \mathrm{(D) \ } 4\qquad \mathrm{(E) \ }7$


 

	\item (\textit{AHSME 1998}) Walter rolls four standard six-sided dice and finds that the product of the numbers of the upper faces is $144$. Which of he following could \textbf{not } be the sum of the upper four faces? 


 $\mathrm{(A) \ }14 \qquad \mathrm{(B) \ }15 \qquad \mathrm{(C) \ }16 \qquad \mathrm{(D) \ }17 \qquad \mathrm{(E) \ }18$


 

	\item (\textit{AHSME 1998}) A parabola has vertex of $(4,-5)$ and has two $x-$intercepts, one positive, and one negative. If this parabola is the graph of $y = ax^2 + bx + c,$ which of $a,b,$ and $c$ must be positive?


 $\mathrm{(A) \ } \text{only}\ a \qquad \mathrm{(B) \ } \text{only}\ b \qquad \mathrm{(C) \ } \text{only}\ c \qquad \mathrm{(D) \ } a\ \text{and}\ b\ \text{only} \qquad \mathrm{(E) \ } \text{none}$


 

	\item (\textit{AHSME 1998}) A regular hexagon and an equilateral triangle have equal areas. What is the ratio of the length of a side of the triangle to the length of a side of the hexagon?


 $\mathrm{(A) \ }\sqrt{3} \qquad \mathrm{(B) \ }2 \qquad \mathrm{(C) \ }\sqrt{6} \qquad \mathrm{(D) \ }3 \qquad \mathrm{(E) \ }6$


 

	\item (\textit{AHSME 1998}) The figure shown is the union of a circle and two semicircles of diameters $a$ and $b$, all of whose centers are collinear. The ratio of the area, of the shaded region to that of the unshaded region is


 
                    \begin{figure}[H]
                        \centering
                        \includegraphics[width=0.5\textwidth]{Problem/AHSME_I_1998_16_Problem_Figure_1.png}
                    \end{figure}
                    \#\#\#px$\mathrm{(A) \ } \sqrt{\frac ab} \qquad \mathrm{(B) \ }\frac ab \qquad \mathrm{(C) \ } \frac{a^2}{b^2} \qquad \mathrm{(D) \ }\frac{a+b}{2b} \qquad \mathrm{(E) \ } \frac{a^2 + 2ab}{b^2 + 2ab}$


 

	\item (\textit{AHSME 1998}) Let $f(x)$ be a function with the two properties:


 
            \begin{enumerate}[label=(\roman*)]
                \item  for any two real numbers $x$ and $y$, $f(x+y) = x + f$, and
\item  $f(0) = 2.$
            \end{enumerate}
            What is the value of $f(1998)?$


 $\mathrm{(A)}\ 0  \qquad\mathrm{(B)}\ 2  \qquad\mathrm{(C)}\ 1996 \qquad\mathrm{(D)}\ 1998  \qquad\mathrm{(E)}\ 2000$


 

	\item (\textit{AHSME 1998}) A right circular cone of volume $A$, a right circular cylinder of volume $M$, and a sphere of volume $C$ all have the same radius, and the common height of the cone and the cylinder is equal to the diameter of the sphere. Then


 $\mathrm{(A) \ } A-M+C = 0 \qquad \mathrm{(B) \ } A+M=C \qquad \mathrm{(C) \ } 2A = M+C \qquad \mathrm{(D) \ }A^2 - M^2 + C^2 = 0 \qquad \mathrm{(E) \ } 2A + 2M = 3C$


 

	\item (\textit{AHSME 1998}) How many triangles have area $10$ and vertices at $(-5,0),(5,0)$ and $(5\cos \theta, 5\sin \theta)$ for some angle $\theta$?


 $\mathrm{(A) \ }0 \qquad \mathrm{(B) \ }2 \qquad \mathrm{(C) \ }4 \qquad \mathrm{(D) \ }6 \qquad \mathrm{(E) \ } 8$


 

	\item (\textit{AHSME 1998}) Three cards, each with a positive integer written on it, are lying face-down on a table. Casey, Stacy, and Tracy are told that


 
            \begin{enumerate}[label=(\roman*)]
                \item  the numbers are all different,
\item  they sum to $13$, and
\item  they are in increasing order, left to right.
            \end{enumerate}
            First, Casey looks at the number on the leftmost card and says, "I don't have enough information to determine the other two numbers."  Then Tracy looks at the number on the rightmost card and says, "I don't have enough information to determine the other two numbers."  Finally, Stacy looks at the number on the middle card and says, "I don't have enough information to determine the other two numbers."  Assume that each person knows that the other two reason perfectly and hears their comments.  What number is on the middle card?


 $\textrm{(A)}\ 2 \qquad \textrm{(B)}\ 3 \qquad \textrm{(C)}\ 4 \qquad \textrm{(D)}\ 5 \qquad \textrm{(E)}\ \text{There is not enough information to determine the number.}$


 

	\item (\textit{AHSME 1998}) In an $h$-meter race, Sunny is exactly $d$ meters ahead of Windy when Sunny finishes the race. The next time they race, Sunny sportingly starts $d$ meters behind Windy, who is at the starting line. Both runners run at the same constant speed as they did in the first race. How many meters ahead is Sunny when Sunny finishes the second race?


 $\mathrm{(A) \ } \frac dh \qquad \mathrm{(B) \ } 0 \qquad \mathrm{(C) \ } \frac {d^2}h \qquad \mathrm{(D) \ } \frac {h^2}d \qquad \mathrm{(E) \ } \frac{d^2}{h-d}$


 

	\item (\textit{AHSME 1998}) What is the value of the expression


 $\frac {1}{\log_{2}100!} + \frac {1}{\log_{3}100!} + \frac {1}{\log_{4}100!} + \cdots + \frac {1}{\log_{100}100!}$

 $\mathrm{(A) \ }0.01 \qquad \mathrm{(B) \ }0.1 \qquad \mathrm{(C) \ }1 \qquad \mathrm{(D) \ }2 \qquad \mathrm{(E) \ } 10$


 

	\item (\textit{AHSME 1998}) The graphs of $x^2 + y^2 = 4 + 12x + 6y$ and $x^2 + y^2 = k + 4x + 12y$ intersect when $k$ satisfies $a \le k \le b$, and for no other values of $k$. Find $b-a$.


 $\mathrm{(A) \ }5 \qquad \mathrm{(B) \ }68 \qquad \mathrm{(C) \ }104 \qquad \mathrm{(D) \ }140 \qquad \mathrm{(E) \ }144$


 

	\item (\textit{AHSME 1998}) Call a $7$-digit telephone number $d_1d_2d_3-d_4d_5d_6d_7$ \textit{memorable} if the prefix sequence $d_1d_2d_3$ is exactly the same as either of the sequences $d_4d_5d_6$ or $d_5d_6d_7$ (possibly both). Assuming that each $d_i$ can be any of the ten decimal digits $0,1,2, \ldots 9$, the number of different memorable telephone numbers is 


 $\mathrm{(A) \ } 19,810 \qquad \mathrm{(B) \ } 19,910 \qquad \mathrm{(C) \ } 19,990 \qquad \mathrm{(D) \ } 20,000 \qquad \mathrm{(E) \ } 20,100$


 

	\item (\textit{AHSME 1998}) A piece of graph paper is folded once so that $(0,2)$ is matched with $(4,0)$, and $(7,3)$ is matched with $(m,n)$. Find $m+n$.


 $\mathrm{(A) \ }6.7 \qquad \mathrm{(B) \ }6.8 \qquad \mathrm{(C) \ }6.9 \qquad \mathrm{(D) \ }7.0 \qquad \mathrm{(E) \ }8.0$


 

	\item (\textit{AHSME 1998}) In quadrilateral $ABCD$, it is given that $\angle A = 120^{\circ}$, angles $B$ and $D$ are right angles, $AB = 13$, and $AD = 46$. Then $AC=$


 $\mathrm{(A) \ } 60 \qquad \mathrm{(B) \ }62 \qquad \mathrm{(C) \ }64 \qquad \mathrm{(D) \ }65 \qquad \mathrm{(E) \ } 72$


 

	\item (\textit{AHSME 1998}) A $9 \times 9 \times 9$ cube is composed of twenty-seven $3 \times 3 \times 3$ cubes. The big cube is ‘tunneled’ as follows: First the six $3 \times 3 \times 3$ cubes which make up the center of each face as well as the center $3 \times 3 \times 3$ cube are removed as shown. Second, each of the twenty remaining $3 \times 3 \times 3$ cubes is diminished in the same way. That is, the center facial unit cubes as well as each center cube are removed. The surface area of the final figure is 


 
                    \begin{figure}[H]
                        \centering
                        \includegraphics[width=0.5\textwidth]{Problem/AHSME_I_1998_27_Problem_Figure_1.png}
                    \end{figure}
                    $\mathrm{(A) \ } 384 \qquad \mathrm{(B) \ } 729 \qquad \mathrm{(C) \ } 864 \qquad \mathrm{(D) \ } 1024 \qquad \mathrm{(E) \ } 1056$


 

	\item (\textit{AHSME 1998}) In triangle $ABC$, angle $C$ is a right angle and $CB > CA$. Point $D$ is located on $\overline{BC}$ so that angle $CAD$ is twice angle $DAB$. If $AC/AD = 2/3$, then $CD/BD = m/n$, where $m$ and $n$ are relatively prime positive integers. Find $m+n$.


 $\mathrm{(A) \ }10 \qquad \mathrm{(B) \ }14 \qquad \mathrm{(C) \ }18 \qquad \mathrm{(D) \ }22 \qquad \mathrm{(E) \ } 26$


 

	\item (\textit{AHSME 1998}) A point $(x,y)$ in the plane is called a lattice point if both $x$ and $y$ are integers. The area of the largest square that contains exactly three lattice points in its interior is closest to


 $\mathrm{(A) \ } 4.0 \qquad \mathrm{(B) \ } 4.2 \qquad \mathrm{(C) \ } 4.5 \qquad \mathrm{(D) \ } 5.0 \qquad \mathrm{(E) \ }  5.6$


 

	\item (\textit{AHSME 1998}) For each positive integer $n$, let 


 $a_n = \frac{(n+9)!}{(n-1)!}$

 Let $k$ denote the smallest positive integer for which the rightmost nonzero digit of $a_k$ is odd. The rightmost nonzero digit of $a_k$ is


 $\mathrm{(A) \ }1 \qquad \mathrm{(B) \ }3 \qquad \mathrm{(C) \ }5 \qquad \mathrm{(D) \ } 7 \qquad \mathrm{(E) \ } 9$


 

\end{enumerate}